% TOP_PROBLEM: {"id": 2617, "title": "Kourovka Notebook Problem 21.108", "category_id": 17, "category": {"id": 17, "name": "group_theory", "display_name": "Group Theory", "description": "Problems about groups, group actions, representations, and related algebraic structures.", "slug": "group-theory", "order_index": 17, "created_at": "2026-07-31T15:26:25.671Z"}, "status": "open", "problem_number": "KOU-21.108", "set_id": 10, "statusQualification": "The omitted inequality defining k is restored from the supplied background: |pi(G)|<=k sigma(G)."}
% FIRST_POSTED: 2026-10-01
% ENTRY_KIND: statement_only
% UnsolvedMath ID 2617; source catalogue code KOU-21.108
% !TeX program = lualatex
% Definition and sources only; this file is not an LLM solution attempt.
\documentclass[11pt,a4paper]{article}
\usepackage[margin=25mm]{geometry}
\usepackage{fontspec}
\setmainfont{lmroman10-regular.otf}[BoldFont=lmroman10-bold.otf,
 ItalicFont=lmroman10-italic.otf,BoldItalicFont=lmroman10-bolditalic.otf]
\usepackage{amsmath,amssymb,mathtools}
\usepackage{unicode-math}
\setmathfont{latinmodern-math.otf}
\AtBeginDocument{\let\setminus\smallsetminus}
\usepackage{xurl,hyperref}
\hypersetup{colorlinks=true,urlcolor=blue}
\setcounter{secnumdepth}{0}
\setlength{\parindent}{0pt}
\setlength{\parskip}{5pt}
\setlength{\emergencystretch}{3em}
\begin{document}
\section*{Kourovka Notebook Problem 21.108}\label{2617}
{\small UnsolvedMath ID 2617 · KOU-21.108 · Group Theory · Kourovka Notebook - New Problems, Issue 21}
\subsection{Definitions and mathematical statement}
For a finite group $G$, let $\operatorname{Irr}(G)$ be its irreducible complex characters, one for each isomorphism class of irreducible complex representation. For $\chi\in\operatorname{Irr}(G)$, its degree is $\chi(1)$ and its kernel is the kernel of a representation affording it; equivalently $\ker\chi=\{g:\chi(g)=\chi(1)\}$. Its codegree is the positive integer
\[
 \operatorname{cod}(\chi)=\frac{[G:\ker\chi]}{\chi(1)}.
\]
Let $\operatorname{Cod}(G)=\{\operatorname{cod}(\chi):\chi\in\operatorname{Irr}(G)\}$ be the set of distinct codegrees.

For a positive integer $m$, let $\pi(m)$ be its set of prime divisors, with $\pi(1)=\varnothing$. Define
\[
 \pi(G)=\pi(|G|),\qquad
 \sigma(G)=\max_{m\in\operatorname{Cod}(G)}|\pi(m)|.
\]
The maximum counts distinct prime divisors, without multiplicity. The principal character has codegree one, so the defining set is nonempty even for the trivial group.

The problem asks whether the universal bound
\[
 |\pi(G)|\le4\sigma(G)
\]
holds for every finite group $G$. This supplies explicitly the inequality referenced by ``the constant $k$'' in the abbreviated catalogue statement. In particular, the question compares all primes occurring in the group order with the largest number that occur together in any single character codegree. No restriction to solvable groups is imposed.
\subsection{Short English statement}
Is the number of primes dividing a finite group’s order at most four times the largest number of primes dividing one of its character codegrees?
\subsection{Sources}
\begin{itemize}
\item[S1] ULAM.AI and the UnsolvedMath Contributors, supplied problems.json, record 2617 (KOU-21.108), Kourovka Notebook - New Problems, Issue 21. Catalogue identity and source transcription; CC BY 4.0.
\url{https://huggingface.co/datasets/ulamai/UnsolvedMath}.
\item[S2] The Kourovka Notebook, 21st edition, version 47 (30 September 2026), new problems, Problem 21.108. Expanded formulation of the supplied catalogue statement.
\url{https://arxiv.org/pdf/1401.0300v47}.
\end{itemize}
\end{document}
